\chapter{Как читать этот курс}

\section{Не FAQ, а курс}

Эта версия текста сознательно написана не как памятка для уже посвящённых и не как
монография, где читатель обязан заранее жить внутри терминологии. Жанр здесь другой:
\emph{курс по архитектуре памяти}. Это значит, что изложение должно делать не только
две вещи — сообщать решения и хранить корректность, — но и третью: показывать,
\emph{почему} решение вообще появилось, на какой вопрос отвечает и чем рискует.

Отсюда три обязательства текста.
\begin{enumerate}
\item Определение появляется раньше, чем начинает незаметно работать.
\item Формула появляется там, где без неё теряется механизм, а не для украшения.
\item Инженерный компромисс называется компромиссом, а не тайком выдаётся за закон природы.
\end{enumerate}

\section{Три уровня утверждений}

В курсе важно не смешивать три рода высказываний.

\begin{definition}[Архитектурное утверждение]
Утверждение архитектурно, если от него зависит сам класс строящейся системы:
например, что память хранится как явный типизированный граф; что retrieval,
внутренняя жизнь и сон реализуются как политики одного обходчика; что факт и
гипотеза различаются эпистемически, а не только текстово.
\end{definition}

\begin{definition}[Инженерное упрощение]
Утверждение инженерно, если оно выбрано как первое рабочее приближение, но в
принципе может быть заменено без смены класса системы: например, жёсткий порог
пересечения вместо полноценной кластерной консолидации, LSH/MinHash вместо более
тонкой структуры сходства, разреженный Inspector вместо тяжёлого explainability-слоя.
\end{definition}

\begin{definition}[Исследовательская ветвь]
Утверждение относится к исследовательской ветви, если оно не нужно для Stage~1,
но обещает новый класс эффектов: клеточные комплексы, динамическая память, резервуарные
вычисления, квантовые ассоциативные схемы.
\end{definition}

\begin{pitfall}[Самая частая методологическая ошибка]
Если инженерное упрощение обсуждается как будто это и есть архитектура, система
становится хрупкой не из-за механики, а из-за языка описания. Разработчик начинает
бояться менять то, что на самом деле было временным. Обратная ошибка не менее вредна:
если архитектурное решение выдаётся за временную деталь, его начинают размывать
«локальными оптимизациями», и система теряет идентичность.
\end{pitfall}

\section{Что именно расширено по сравнению с базовым дублем}

В этом курсе есть места, где добавлена не только педагогика, но и техника.
Это сделано открыто.

\begin{enumerate}
\item Retrieval расписан не только качественно, но и как оператор на графе: с
матрицей переходов, рестартом, температурой, hitting time и спектральными
ограничениями.
\item Эпистемика развернута в виде каналов свидетельств, а не только одного
скалярного поля уверенности.
\item Экспериментальный раздел усилен явной диагностикой failure modes.
\item Добавлены сквозные разборы, где память прослеживается от записи события до
ревизии производных узлов.
\end{enumerate}

\begin{intuition}[Почему это не «раздуть текст»]
Когда архитектура сложна, краткость бывает двух разных сортов. Бывает краткость
хорошая: выкинули повтор и оставили идею. А бывает краткость опасная: скрыли
промежуточные мосты, и тогда текст выглядит компактным, но держится на том, что
читатель сам за автора достроит половину механизма. Расширенный курс борется именно
со вторым типом краткости.
\end{intuition}

\section{Три прохода чтения}

Для такого текста полезно иметь не один, а три режима чтения.

\paragraph{Первый проход — карта.} Нужно понять, что система вообще считает памятью,
какие в ней есть процессы и где граница между графом, логом, рантаймом и языковой
моделью.

\paragraph{Второй проход — механика.} Здесь уже важны формулы, типы узлов, поля,
операторы консолидации, режимы обходчика и пропагация уверенности.

\paragraph{Третий проход — споры и открытые ветви.} На этом уровне читаются не
столько ответы, сколько развилки: где исходное решение твёрдое, где оно временное,
что можно заменить, а что нельзя, не разрушив замысел.

\section{Главный критерий хорошего текста о памяти}

Хороший текст об архитектуре памяти должен выдерживать двоякую проверку.

\begin{principle}[Два экзамена]
Первый экзамен — \emph{технический}: после чтения должно быть возможно собрать Stage~1,
придумать тесты, найти диагностические переменные и не потерять roadmap. Второй
экзамен — \emph{педагогический}: должно быть понятно, почему эта архитектура вообще
выглядит именно так, а не как список случайно соседствующих трюков.
\end{principle}

Этим двум экзаменам и подчинён весь дальнейший порядок курса.
